\section*{Chapter \ref{ch:g9:gravitation} --- Gravitation}

\begin{solution}{exo:g9:gravitation:1}
$F = G \, m_A m_B / d^2$, with masses in \omterm{def:g3:measuring-mass:units}{kilograms}, $d$ in
\omterm{def:g2:measuring-length:units}{metres} between centers, $F$ in newtons and
$G = \qty{6.67e-11}{N.m^2/kg^2}$. In words: the pull grows
with each \omterm{def:g3:measuring-mass:mass}{mass} and dies with the square of the distance ---
and it acts equally on both partners, in opposite directions.
\end{solution}

\begin{solution}{exo:g9:gravitation:2}
With exactly \qty{1}{N} --- the law's pulls always come in
equal, opposite pairs. The apple visibly accelerates because
its \omterm{def:g3:measuring-mass:mass}{mass} is tiny; the same \qty{1}{N} applied to
\qty{5.97e24}{kg} of \omterm{def:g4:solar-system:planet}{planet} stirs it immeasurably little.
\end{solution}

\begin{solution}{exo:g9:gravitation:3}
At $2d$: a quarter of the \omterm{def:g2:pushes-and-pulls:force}{force} at $d$ --- the square of two.
At $10d$: a hundredth.
\end{solution}

\begin{solution}{exo:g9:gravitation:4}
$F = \num{6.67e-11} \times 80 \times 80 \div 0.25 =
\qty{1.7e-6}{N}$ --- under two millionths of a newton:
whatever draws dance partners together, it is not
\omterm{def:g9:gravitation:gravitation}{gravitation}.
\end{solution}

\begin{solution}{exo:g9:gravitation:5}
The rate at which the ball falls, and the rate at which the
round Earth's surface curves away beneath it. An \omterm{def:g6:sun-earth-moon:axisorbit}{orbit}:
falling and missing.
\end{solution}

\begin{solution}{exo:g9:gravitation:6}
At the station's altitude, gravity keeps nearly its full
surface strength --- it is what holds the station on its
circle. Station and crew are both \emph{falling}, together,
around the Earth: falling companions press on nothing, and
floating is the feeling of shared free fall.
\end{solution}

\begin{solution}{exo:g9:gravitation:7}
Because the heavens deal in monstrous masses --- \omterm{def:g4:solar-system:planet}{planets} and
\omterm{def:g4:solar-system:star}{stars} multiply the feeble $G$ into mighty \omterm{def:g2:pushes-and-pulls:force}{forces} --- and
because \omterm{def:g9:gravitation:gravitation}{gravitation} is never canceled: it only \omterm{def:g1:magnets:magnet}{attracts}, so
every added \omterm{def:g3:measuring-mass:mass}{mass} deepens the pull.
\end{solution}

\begin{solution}{exo:g9:gravitation:8}
The \omterm{def:g2:moon-in-the-sky:moon}{Moon}'s pull heaps the seas both toward it (nearest water,
pulled hardest) and away from it (farthest water, pulled
least); the spinning Earth carries each coast under both
heaps daily --- two tides. Greatest at new and \omterm{ex:g2:moon-in-the-sky:shapes}{full moon},
when the Sun's hand aligns with the \omterm{def:g2:moon-in-the-sky:moon}{Moon}'s: spring tides.
\end{solution}

\begin{solution}{exo:g9:gravitation:9}
$F = \num{6.67e-11} \times (\num{7.3e22} \times
\num{5.97e24}) \div (\num{3.8e8})^2 \approx \qty{2.0e20}{N}$
--- a $2$ followed by twenty zeros of newtons holding the
\omterm{def:g2:moon-in-the-sky:moon}{Moon} on its lease.
\end{solution}

\begin{solution}{exo:g9:gravitation:10}
Two dials turn against each other: the \omterm{def:g2:moon-in-the-sky:moon}{Moon}'s \omterm{def:g3:measuring-mass:mass}{mass}, about
eighty times smaller, cuts the pull eightyfold; but its
radius, nearly four times smaller, \emph{boosts} the surface
pull by four squared --- about fourteenfold. Eighty down,
fourteen up: near a sixth remains.
\end{solution}

\begin{solution}{exo:g9:gravitation:11}
The inverse square thins the Sun's grip with distance, and a
weaker grip can hold only a slower fall-around: the outer
\omterm{def:g4:solar-system:planet}{planets} are leashed loosely and amble, the inner ones are
gripped hard and race --- Mercury's sprint and Neptune's
164-year stroll are the same law read at two distances.
\end{solution}

\begin{solution}{exo:g9:gravitation:12}
Its period must be exactly one day, circling eastward --- the
Earth's own spin direction. Then satellite and antenna turn in
perfect step: the ``hanging'' dish is watching a cannonball
that falls around the \omterm{def:g4:solar-system:planet}{planet} at precisely the rate the \omterm{def:g4:solar-system:planet}{planet}
rotates beneath it --- motionless only in the spinning
observer's eyes.
\end{solution}

\begin{solution}{pb:g9:gravitation:1}
\textbf{1.} $F = \num{6.67e-11} \times (0.1 \times
\num{5.97e24}) \div (\num{6.4e6})^2 \approx \qty{0.97}{N}$
--- call it one newton: the orchard tug, computed from the
universe's constant.
\textbf{2.} \qty{0.97}{N} likewise --- mutuality: the law's
pulls come only in equal, opposite pairs.
\textbf{3.} Doubling the center distance quarters the pull:
about \qty{0.24}{N}.
\textbf{4.} Every crumb of the \omterm{def:g4:solar-system:planet}{planet} pulls the apple ---
mountains near, core deep, antipodes far --- and the sum of
all those pulls works out exactly as if the whole \omterm{def:g3:measuring-mass:mass}{mass} sat at
the center: a theorem, not an assumption, and its author
delayed publishing for years until he could prove it.
\textbf{5.} About $2000 \times 9.8 \approx \qty{2.0e4}{N}$
--- twenty thousand newtons (the law with $d$ one Earth
radius gives the same).
\textbf{6.} At two radii, the inverse square leaves a
quarter: about \qty{4900}{N} per its \qty{19600}{N} surface
\omterm{def:g4:weight-and-mass:weight}{weight}.
\textbf{7.} Yes --- a quarter of full gravity is no
weightless void. The satellite spends that pull on curving
its path: it falls around the Earth, forever missing, exactly
as the mountain cannon taught.
\textbf{8.} ``Fired fast enough, a falling cannonball's drop
matches the Earth's curve, and it circles forever ---
falling without landing. Every satellite is that cannonball,
launched sideways fast enough to keep missing the ground.''
\textbf{9.} $\num{3.8e8} \div \num{6.4e6} \approx 59.4$
Earth radii.
\textbf{10.} About $1 \div 59.4^2 \approx 1/3500$ of surface
strength.
\textbf{11.} Gentle because gravity at sixty radii is
diluted three-and-a-half-thousandfold --- the \omterm{def:g2:measuring-length:units}{five-metre}
first-second drop shrinks to about a millimetre. Endless
because that millimetre of fall exactly matches the
millimetre by which the Earth's surface curves away beneath
the \omterm{def:g2:moon-in-the-sky:moon}{Moon}'s sideways glide: the fall never gains.
\textbf{12.} For example: ``Apple, satellite and \omterm{def:g2:moon-in-the-sky:moon}{Moon} obey
one sentence: every \omterm{def:g3:measuring-mass:mass}{mass} falls toward every other, by
$G m_A m_B / d^2$ --- the apple lands, the satellite and the
\omterm{def:g2:moon-in-the-sky:moon}{Moon} keep missing, and the difference is only sideways
\omterm{def:g7:motion-average-speed:formula}{speed}.''
\end{solution}
